% Preamble requirements:
% \usepackage{amsmath,amssymb,booktabs,array,longtable,enumitem,tikz}
% \usetikzlibrary{arrows.meta,positioning,shapes.geometric}
%
% Bibliography files used by this section:
%   the proposal references.bib
%   doc/related_work_additions.bib
%   doc/dataset_preprocessing_additions.bib

\providecommand{\datasetfield}[1]{\texttt{\detokenize{#1}}}

\section{Dataset Construction and Preprocessing}
\label{sec:dataset_preprocessing}

This section describes how the proprietary packet capture was converted into
the two aligned data views required by the proposed hierarchical detector. The
description distinguishes among (i) fields retained in the raw CSV files,
(ii) predictors supplied to the Stage~1 model, and (iii) identifiers retained
only for constructing Stage~2 context. This distinction is important because a
field can be useful for record linkage without being an appropriate learning
feature. It also makes explicit which operations are fitted on training data
and which operations are deterministic transformations of the capture.
Careful documentation of provenance, labels, transformations, and partitions
is essential for credible NIDS evaluation
\cite{Shiravi2012DatasetGeneration,Ring2019Datasets,Goldschmidt2025Datasets,
Pinto2025DatasetReview}.

\subsection{Data Provenance and Scope}
\label{subsec:data_provenance}

The experiments use the packet capture
\datasetfield{ar002_et12_20260511_001.pcap}, whose file size is approximately
$8.62\,\mathrm{GB}$. The trace was supplied by a company colleague after being
downloaded from the company network with an organizational traffic-collection
tool. The thesis author therefore starts from an already acquired PCAP rather
than generating traffic in a laboratory testbed. Unlike CICIDS-style datasets,
in which benign profiles and attacks are deliberately executed under a
controlled schedule \cite{Sharafaldin2018CICIDS}, this capture represents
naturally occurring enterprise traffic as observed by the collection point.
It may consequently provide greater environmental realism, but its traffic
mix, host population, visibility, and attack coverage are specific to one
organization and one capture period.

The dataset is proprietary and is not presented as a public benchmark. The raw
PCAP and CSV files contain network identifiers, including IP addresses and
ports, that can be sensitive. Accordingly, this thesis reports only aggregate
statistics and does not print individual addresses. Any artifact released
outside the authorized environment should replace host identifiers with
consistent pseudonyms, remove unnecessary endpoint fields, and be reviewed
against the company's data-governance policy. Consistent pseudonyms are
preferable to deleting all identifiers because Stage~2 requires equality tests
between host identifiers; the mapping itself need not be released. The current
extraction code retains raw addresses in local metadata, so pseudonymization is
a release requirement rather than a claimed property of the present CSV files.

\subsection{PCAP-to-CSV Extraction}
\label{subsec:pcap_extraction}

The PCAP is replayed through Suricata, an open-source network threat detection
engine \cite{Roesch1999,SuricataGuide2026}. A custom C output plugin located at
\datasetfield{suricata-plugin/custom-out.c} registers both a packet logger and
a flow logger. Suricata exposes low-level packet and flow logging interfaces to
custom output modules; the plugin uses these interfaces to emit two CSV files
in one processing pass. This design avoids reconstructing packet-to-flow
membership in a separate tool and preserves a shared bidirectional
\datasetfield{flow_id} across both views.

Figure~\ref{fig:data_construction_pipeline} summarizes the data-construction
path. Packet rows preserve within-flow order and microsecond timestamps, while
the flow logger maintains online aggregates and writes one statistical record
per completed flow. The flow schema follows the broad feature families used by
CICFlowMeter/CICIDS-style data, while adding explicit indicators for whether
the initial forward and backward TCP window values are available
\cite{Sharafaldin2018CICIDS}.

\begin{figure}[t]
    \centering
    \begin{tikzpicture}[
        node distance=7mm and 8mm,
        block/.style={draw, rectangle, rounded corners=2pt, align=center,
                      minimum height=9mm, text width=29mm, font=\small},
        smallblock/.style={draw, rectangle, rounded corners=2pt, align=center,
                           minimum height=9mm, text width=27mm, font=\small},
        arrow/.style={-{Latex[length=2mm]}, thick}
    ]
        \node[block] (pcap) {Enterprise PCAP\\$8.62\,\mathrm{GB}$};
        \node[block, right=of pcap] (suricata) {Suricata + custom\\output plugin};
        \node[smallblock, above right=5mm and 9mm of suricata] (packet) {Packet CSV\\$11{,}750{,}727\times40$};
        \node[smallblock, below right=5mm and 9mm of suricata] (flow) {Flow CSV\\$421{,}898\times77$};
        \node[block, right=9mm of packet] (align) {Integrity checks and\\\datasetfield{flow_id} alignment};
        \node[block, below=of align] (split) {Chronological\\flow-level split};
        \node[block, below=of split] (tensor) {Train-fitted transforms\\and model tensors};

        \draw[arrow] (pcap) -- (suricata);
        \draw[arrow] (suricata) -- (packet);
        \draw[arrow] (suricata) -- (flow);
        \draw[arrow] (packet) -- (align);
        \draw[arrow] (flow) -- (align);
        \draw[arrow] (align) -- (split);
        \draw[arrow] (split) -- (tensor);
    \end{tikzpicture}
    \caption{Construction and preprocessing path from the proprietary PCAP to
    Stage~1 tensors and Stage~2 metadata.}
    \label{fig:data_construction_pipeline}
\end{figure}

No raw application payload bytes are supplied to either learning stage. The
packet table does contain \datasetfield{payload_len}, which is a scalar length
derived from packet metadata, but it does not contain payload content. This
wording is deliberate: the detector is payload-content independent, not
entirely free of payload-related metadata. Alert signatures, rule identifiers,
alert messages, and alert counts are also excluded from the predictors.

\subsection{Alert-Derived Label Construction}
\label{subsec:label_construction}

The custom logger derives binary labels from Suricata's alert state. Let
$a(p_i)$ denote the number of Suricata alerts attached to packet $p_i$. The
packet label is
\begin{equation}
    y_i^{\mathrm{pkt}}
    = \mathbb{I}\!\left[a(p_i)>0\right],
    \label{eq:packet_alert_label}
\end{equation}
where $\mathbb{I}[\cdot]$ is the indicator function. A flow is labelled
positive when at least one of its packets triggers an alert:
\begin{equation}
    y_f
    = \max_{p_i\in\mathcal{P}_f} y_i^{\mathrm{pkt}},
    \label{eq:flow_alert_label}
\end{equation}
where $\mathcal{P}_f$ is the set of packets assigned to flow $f$. Thus,
\datasetfield{packet_label}=1 indicates an alert-triggering packet and
\datasetfield{label}=1 indicates an alert-associated flow.

These are programmatic or \emph{weak} labels rather than independent expert
ground truth \cite{Ratner2017Snorkel}. In particular, label~0 means that the
configured Suricata rules did not generate an alert; it does not prove that the
traffic is benign. Conversely, a positive label inherits any false positives
of the deployed ruleset. The experiments therefore measure how well the model
learns and contextualizes Suricata-alert-associated behavior, rather than
establishing guaranteed detection of every intrusion or zero-day attack.
Potential label errors can affect both training and test estimates, and this
limitation must be retained when interpreting the results
\cite{Northcutt2021ConfidentLearning}.

To prevent direct target leakage, neither \datasetfield{packet_label} nor
\datasetfield{label} is included in a feature vector. Suricata alert metadata
is likewise absent from the model input. For reproducibility, the final
experiment manifest should record the Suricata version, rule-set name and
version, rule update date, command-line/configuration checksum, custom-plugin
commit, and PCAP checksum. These items determine the weak-labeling function and
cannot be reconstructed from the CSV dimensions alone.

\subsection{Raw Dataset Composition}
\label{subsec:raw_composition}

The extraction produced the two aligned tables summarized in
Table~\ref{tab:raw_dataset_summary}. The packet table contains
$11{,}750{,}727$ rows and 40 columns; the flow table contains $421{,}898$ rows
and 77 columns. Dividing the two row counts gives an average of approximately
$27.85$ logged packets per flow before sequence truncation. This average does
not imply that flow lengths are uniform; network packet counts are typically
heavy-tailed, which motivates an explicit maximum sequence length.

\begin{table}[t]
    \centering
    \caption{Raw extracted dataset summary.}
    \label{tab:raw_dataset_summary}
    \begin{tabular}{@{}lrrp{5.7cm}@{}}
        \toprule
        Data view & Rows & Columns & Primary role \\
        \midrule
        Packet CSV & 11,750,727 & 40 & Ordered packet-header, protocol,
        direction, and inter-arrival observations within flows \\
        Flow CSV & 421,898 & 77 & One flow-level target, aggregate
        statistics, and context-construction metadata per flow \\
        \bottomrule
    \end{tabular}
\end{table}

The binary flow-label distribution is shown in
Table~\ref{tab:flow_label_distribution}. Of the $421{,}898$ flow records,
$405{,}254$ are labelled 0 and $16{,}644$ are labelled 1. The positive
prevalence and imbalance ratio are therefore
\begin{align}
    \pi_1 &= \frac{16{,}644}{421{,}898} = 0.03945
             \quad (3.945\%), \\
    \mathrm{IR} &= \frac{405{,}254}{16{,}644} = 24.35:1.
    \label{eq:dataset_imbalance}
\end{align}

\begin{table}[t]
    \centering
    \caption{Flow-level alert-derived label distribution.}
    \label{tab:flow_label_distribution}
    \begin{tabular}{@{}llrr@{}}
        \toprule
        Label & Interpretation & Count & Proportion \\
        \midrule
        0 & No Suricata alert observed in the flow & 405,254 & 96.055\% \\
        1 & At least one alert-triggering packet & 16,644 & 3.945\% \\
        \midrule
        Total & & 421,898 & 100.000\% \\
        \bottomrule
    \end{tabular}
\end{table}

An always-negative classifier would consequently obtain approximately
$96.055\%$ accuracy while detecting no positive flows. Accuracy and weighted
averages are therefore secondary metrics for this dataset. The evaluation
focuses on macro-precision, macro-recall, macro-F1, ROC-AUC, PR-AUC, the
confusion matrix, and the false-positive rate. PR-AUC is particularly useful
when the positive class is rare \cite{SaitoRehmsmeier2015}. The large benign
population also makes the false-positive rate operationally important: a small
percentage of $405{,}254$ negative flows can still produce a substantial alert
volume.

\subsection{Raw Packet and Flow Schemas}
\label{subsec:raw_schemas}

Tables~\ref{tab:packet_schema} and~\ref{tab:flow_schema} group all raw columns
by function. The grouping is descriptive; it does not imply that every raw
column is used as a predictor.

\begin{longtable}{@{}>{\raggedright\arraybackslash}p{0.20\textwidth}
                         >{\raggedright\arraybackslash}p{0.75\textwidth}@{}}
    \caption{Complete packet CSV schema grouped by function.}
    \label{tab:packet_schema} \\
    \toprule
    Group & Raw columns \\
    \midrule
    \endfirsthead
    \multicolumn{2}{c}{\tablename\ \thetable\ -- continued from previous page} \\
    \toprule
    Group & Raw columns \\
    \midrule
    \endhead
    \midrule
    \multicolumn{2}{r}{Continued on next page} \\
    \endfoot
    \bottomrule
    \endlastfoot
    Record identity and time &
    \datasetfield{record_type}, \datasetfield{packet_id},
    \datasetfield{timestamp_us}, \datasetfield{ts_sec},
    \datasetfield{ts_usec}, \datasetfield{flow_id} \\
    Endpoints and protocol &
    \datasetfield{direction}, \datasetfield{ip_version},
    \datasetfield{src_ip}, \datasetfield{dst_ip},
    \datasetfield{src_port}, \datasetfield{dst_port},
    \datasetfield{protocol}, \datasetfield{l7_proto} \\
    Length and header metadata &
    \datasetfield{pkt_len}, \datasetfield{ip_len},
    \datasetfield{payload_len}, \datasetfield{ttl_or_hop_limit},
    \datasetfield{tcp_window}, \datasetfield{tcp_header_len},
    \datasetfield{udp_len} \\
    Transport indicators &
    \datasetfield{has_tcp}, \datasetfield{has_udp},
    \datasetfield{has_icmp} \\
    TCP state and flags &
    \datasetfield{tcp_seq}, \datasetfield{tcp_ack},
    \datasetfield{tcp_flags}, \datasetfield{tcp_fin},
    \datasetfield{tcp_syn}, \datasetfield{tcp_rst},
    \datasetfield{tcp_psh}, \datasetfield{tcp_ack_flag},
    \datasetfield{tcp_urg}, \datasetfield{tcp_ece},
    \datasetfield{tcp_cwr} \\
    ICMP metadata &
    \datasetfield{icmp_type}, \datasetfield{icmp_code} \\
    Inter-arrival time &
    \datasetfield{flow_iat_us}, \datasetfield{direction_iat_us} \\
    Target &
    \datasetfield{packet_label} \\
\end{longtable}

The two packet timing fields have different semantics.
\datasetfield{flow_iat_us} is the elapsed time from the previous packet in the
same bidirectional flow, whereas \datasetfield{direction_iat_us} is measured
from the previous packet in the same direction. Both use microseconds and are
zero for the first applicable observation.

\begin{longtable}{@{}>{\raggedright\arraybackslash}p{0.20\textwidth}
                         >{\raggedright\arraybackslash}p{0.75\textwidth}@{}}
    \caption{Complete flow CSV schema grouped by function.}
    \label{tab:flow_schema} \\
    \toprule
    Group & Raw columns \\
    \midrule
    \endfirsthead
    \multicolumn{2}{c}{\tablename\ \thetable\ -- continued from previous page} \\
    \toprule
    Group & Raw columns \\
    \midrule
    \endhead
    \midrule
    \multicolumn{2}{r}{Continued on next page} \\
    \endfoot
    \bottomrule
    \endlastfoot
    Identity and time &
    \datasetfield{flow_id}, \datasetfield{flow_start_timestamp_us},
    \datasetfield{flow_end_timestamp_us} \\
    Endpoints and protocol &
    \datasetfield{source_ip}, \datasetfield{source_port},
    \datasetfield{destination_ip}, \datasetfield{destination_port},
    \datasetfield{protocol} \\
    Directional volume &
    \datasetfield{total_fwd_packets}, \datasetfield{total_backward_packets},
    \datasetfield{total_length_of_fwd_packets},
    \datasetfield{total_length_of_bwd_packets} \\
    Directional length statistics &
    \datasetfield{fwd_packet_length_max},
    \datasetfield{fwd_packet_length_min},
    \datasetfield{fwd_packet_length_mean},
    \datasetfield{fwd_packet_length_std},
    \datasetfield{bwd_packet_length_max},
    \datasetfield{bwd_packet_length_min},
    \datasetfield{bwd_packet_length_mean},
    \datasetfield{bwd_packet_length_std} \\
    Overall rate and timing &
    \datasetfield{flow_duration}, \datasetfield{flow_bytes_per_s},
    \datasetfield{flow_packets_per_s}, \datasetfield{flow_iat_mean},
    \datasetfield{flow_iat_std}, \datasetfield{flow_iat_max},
    \datasetfield{flow_iat_min} \\
    Directional inter-arrival time &
    \datasetfield{fwd_iat_total}, \datasetfield{fwd_iat_mean},
    \datasetfield{fwd_iat_std}, \datasetfield{fwd_iat_max},
    \datasetfield{fwd_iat_min}, \datasetfield{bwd_iat_total},
    \datasetfield{bwd_iat_mean}, \datasetfield{bwd_iat_std},
    \datasetfield{bwd_iat_max}, \datasetfield{bwd_iat_min} \\
    Directional flags, headers, and rates &
    \datasetfield{fwd_psh_flags}, \datasetfield{bwd_psh_flags},
    \datasetfield{fwd_urg_flags}, \datasetfield{bwd_urg_flags},
    \datasetfield{fwd_header_length}, \datasetfield{bwd_header_length},
    \datasetfield{fwd_packets_per_s}, \datasetfield{bwd_packets_per_s} \\
    Overall packet-length distribution &
    \datasetfield{min_packet_length}, \datasetfield{max_packet_length},
    \datasetfield{packet_length_mean}, \datasetfield{packet_length_std},
    \datasetfield{packet_length_variance} \\
    TCP flag counts &
    \datasetfield{fin_flag_count}, \datasetfield{syn_flag_count},
    \datasetfield{rst_flag_count}, \datasetfield{psh_flag_count},
    \datasetfield{ack_flag_count}, \datasetfield{urg_flag_count},
    \datasetfield{cwe_flag_count}, \datasetfield{ece_flag_count} \\
    Ratios, segments, and initial windows &
    \datasetfield{down_up_ratio}, \datasetfield{average_packet_size},
    \datasetfield{avg_fwd_segment_size},
    \datasetfield{avg_bwd_segment_size},
    \datasetfield{has_init_win_bytes_forward},
    \datasetfield{init_win_bytes_forward},
    \datasetfield{has_init_win_bytes_backward},
    \datasetfield{init_win_bytes_backward},
    \datasetfield{act_data_pkt_fwd}, \datasetfield{min_seg_size_forward} \\
    Active and idle periods &
    \datasetfield{active_mean}, \datasetfield{active_std},
    \datasetfield{active_max}, \datasetfield{active_min},
    \datasetfield{idle_mean}, \datasetfield{idle_std},
    \datasetfield{idle_max}, \datasetfield{idle_min} \\
    Target & \datasetfield{label} \\
\end{longtable}

For a non-TCP packet, TCP-specific values are represented by zero and the
\datasetfield{has_tcp} indicator disambiguates ``not applicable'' from an
observed zero. The same principle is used for UDP and ICMP presence. For flow
records, \datasetfield{has_init_win_bytes_forward} and
\datasetfield{has_init_win_bytes_backward} indicate whether the corresponding
initial-window measurements are available; unavailable window values are
stored as zero.

\subsection{Integrity Checks and Packet--Flow Alignment}
\label{subsec:data_integrity}

The preprocessing code first normalizes column names and removes
\datasetfield{record_type}, which is a logging descriptor rather than a
predictor. Positive and negative infinity are converted to missing values.
Rows with a missing or zero \datasetfield{flow_id} are removed. Flow rows must
also have a nonzero, nonmissing \datasetfield{flow_start_timestamp_us}, and
packet rows must have a nonzero, nonmissing \datasetfield{timestamp_us}.

The flow table is required to contain at most one effective row per
\datasetfield{flow_id}. If duplicate flow identifiers occur, records are
stably ordered by flow identifier and start time and the first row is retained;
the duplicate count is logged. The pipeline then computes the intersection of
flow identifiers found in the packet and flow tables and removes unmatched
records from both sides. Finally, packet records are sorted by
\datasetfield{flow_id} and \datasetfield{timestamp_us}. These checks establish
the one-to-many relation
\begin{equation}
    f \longleftrightarrow
    \left\{p_{f,1},p_{f,2},\ldots,p_{f,n_f}\right\},
    \qquad
    t_{f,1}\leq t_{f,2}\leq\cdots\leq t_{f,n_f},
    \label{eq:packet_flow_alignment}
\end{equation}
before any split or feature transformation is performed. The counts in
Table~\ref{tab:raw_dataset_summary} describe the extracted CSV files; the saved
split report is the authoritative source for post-cleaning counts in a
particular run.

\subsection{Feature Selection and Leakage Control}
\label{subsec:feature_selection}

The model input is configuration-controlled rather than obtained by passing all
raw columns to the network. Table~\ref{tab:stage1_packet_features} lists the 23
raw packet predictors used in the principal Stage~1 configuration before
one-hot expansion. Numerical fields describe packet and header sizes; categorical
fields describe protocol states; and binary fields retain packet direction,
transport presence, and individual TCP flags.

\begin{table}[t]
    \centering
    \caption{Packet predictors used by the principal Stage~1 configuration.}
    \label{tab:stage1_packet_features}
    \small
    \begin{tabular}{@{}p{2.5cm}p{11.4cm}@{}}
        \toprule
        Type & Selected raw columns \\
        \midrule
        Numerical (7) &
        \datasetfield{pkt_len}, \datasetfield{ip_len},
        \datasetfield{payload_len}, \datasetfield{ttl_or_hop_limit},
        \datasetfield{tcp_window}, \datasetfield{tcp_header_len},
        \datasetfield{udp_len} \\
        Categorical (4) &
        \datasetfield{ip_version}, \datasetfield{protocol},
        \datasetfield{icmp_type}, \datasetfield{icmp_code} \\
        Binary (12) &
        \datasetfield{direction}, \datasetfield{has_tcp},
        \datasetfield{has_udp}, \datasetfield{has_icmp},
        \datasetfield{tcp_fin}, \datasetfield{tcp_syn},
        \datasetfield{tcp_rst}, \datasetfield{tcp_psh},
        \datasetfield{tcp_ack_flag}, \datasetfield{tcp_urg},
        \datasetfield{tcp_ece}, \datasetfield{tcp_cwr} \\
        \bottomrule
    \end{tabular}
\end{table}

Raw IP addresses, ports, packet/flow identifiers, absolute timestamps, TCP
sequence and acknowledgment numbers, the combined TCP flag integer,
\datasetfield{l7_proto}, and both target columns are excluded from the Stage~1
predictor matrix. This avoids memorizing site-specific addresses and reduces
the risk that identifiers or absolute capture position become shortcuts for
the target. The principal time tensor is instead derived from
\datasetfield{flow_iat_us}, preserving relative packet timing without exposing
the absolute capture timestamp as an ordinary feature.

For Stage~1 configurations with statistical fusion, a configured subset of the
flow measurements in Table~\ref{tab:flow_schema} is transformed separately.
The principal fusion configuration includes duration, directional packet and
byte counts, packet-length statistics, rates, inter-arrival statistics, flag
counts, header lengths, segment/window measurements, and selected active-period
statistics. The two initial-window availability indicators are treated as
binary features. Endpoint identifiers, endpoint ports, absolute flow times,
and \datasetfield{label} are not part of this statistical vector. The
packet-only ablation disables this flow branch entirely; therefore, raw schema
availability must not be confused with feature use in every experiment.

The fields \datasetfield{flow_start_timestamp_us}, \datasetfield{source_ip},
and \datasetfield{destination_ip} are exported alongside each Stage~1
embedding for Stage~2. They are relationship metadata, not Stage~1 predictors.
Stage~2 compares host identifiers to construct source-host,
destination-host, or endpoint contexts and uses start time to enforce temporal
order. It does not one-hot encode raw addresses as semantic model features and
does not retrieve labels into a context window. This separation limits
identity memorization while preserving the relations required by the proposed
inter-flow model.

\subsection{Train-Fitted Feature Transformations}
\label{subsec:feature_transformations}

All data-dependent transformations are fitted on the training partition only
and then reused without refitting on validation and test data. This ordering is
necessary because fitting clipping bounds, scales, or category vocabularies on
the complete dataset would leak information about future partitions
\cite{Arp2022DosDonts}. Packet and flow numerical features use the same type of
pipeline but have separate fitted parameters.

\paragraph{Numerical features.}
Values are converted to numeric form and non-finite values are mapped to
missing. Selected nonnegative, heavy-tailed features are transformed as
\begin{equation}
    \widetilde{x}=\log(1+\max(x,0)).
    \label{eq:log1p_preprocessing}
\end{equation}
The configured log transform covers packet lengths, payload length, TCP window,
header length, and UDP length, together with many count, duration, byte, rate,
and inter-arrival flow statistics. Missing values are then filled with zero.

For each numerical feature $j$, the training partition provides lower and
upper quantiles $q_{0.001,j}^{\mathrm{tr}}$ and
$q_{0.999,j}^{\mathrm{tr}}$. Values in every split are winsorized by
\begin{equation}
    x_{ij}^{\mathrm{clip}}
    = \min\!\left(
        \max\!\left(\widetilde{x}_{ij},q_{0.001,j}^{\mathrm{tr}}\right),
        q_{0.999,j}^{\mathrm{tr}}
      \right).
    \label{eq:quantile_clipping}
\end{equation}
The principal configuration subsequently applies robust scaling:
\begin{equation}
    x_{ij}^{\mathrm{scaled}}
    = \frac{x_{ij}^{\mathrm{clip}}-
             \operatorname{median}_{\mathrm{tr}}(x_j)}
            {\operatorname{IQR}_{\mathrm{tr}}(x_j)},
    \label{eq:robust_scaling}
\end{equation}
where the denominator is the training interquartile range. Log compression,
tail clipping, and robust scaling together reduce the influence of extreme
traffic bursts without using test statistics.

\paragraph{Categorical features.}
Missing categorical values are represented by the explicit token
\datasetfield{MISSING}. A one-hot encoder is fitted on training categories,
and unseen validation or test categories are ignored rather than causing an
error or expanding the model input dimension.

\paragraph{Binary features.}
Binary fields are converted to numeric form, non-finite and missing entries are
set to zero, positive values are mapped to one, and all other values are mapped
to zero. Binary variables are not one-hot encoded, log-transformed, clipped, or
scaled.

\subsection{Packet-Sequence Construction}
\label{subsec:sequence_construction}

For each aligned flow $f$, packets are ordered by
\datasetfield{timestamp_us}. The principal configuration uses a maximum
sequence length $L=128$ and the \datasetfield{head} selection strategy. If
$n_f>L$, the first 128 packets are retained; if $n_f<L$, the sequence is
right-padded with zeros. A Boolean mask identifies valid positions so that
padding does not contribute to attention or pooling. The resulting arrays have
the forms
\begin{equation}
    \mathbf{X}\in\mathbb{R}^{N\times L\times d_p},\qquad
    \mathbf{T}\in\mathbb{R}^{N\times L},\qquad
    \mathbf{M}\in\{0,1\}^{N\times L},
    \label{eq:stage1_tensor_shapes}
\end{equation}
where $d_p$ is the transformed packet-feature dimension after one-hot
encoding.

The timing tensor is computed from the nonnegative within-flow inter-arrival
time:
\begin{equation}
    T_{f,i}=\log\!\left(1+max(Delta t_{f,i},0)\right),
    \qquad
    \Delta t_{f,i}=\datasetfield{flow_iat_us}(p_{f,i}).
    \label{eq:time_tensor}
\end{equation}
It is supplied separately to the time-aware encoding rather than concatenated
with absolute timestamps.

Two alternative truncation strategies are retained for ablation experiments.
\datasetfield{head_tail} keeps the first $\lfloor L/2\rfloor$ and last
$L-\lfloor L/2\rfloor$ packets, whereas \datasetfield{random} draws $L$
packets without replacement using a deterministic seed derived from the global
seed and \datasetfield{flow_id}. Selected random packets are restored to
chronological order. Sequence lengths $8,16,32,64,128,$ and $256$ are evaluated
separately; these ablations change the retained evidence and computational
cost, but not the underlying flow split.

\subsection{Chronological Train--Validation--Test Partition}
\label{subsec:chronological_partition}

Splitting is performed at the flow level after sorting unique flows by
\datasetfield{flow_start_timestamp_us} with a stable ordering. The principal
configuration assigns the earliest 70\% of flows to training, the next 10\%
to validation, and the latest 20\% to testing. It sets
\datasetfield{stratify=false}; therefore, the reported main split uses the
nominal chronological boundaries and does not move examples to equalize class
ratios. Assuming the $421{,}898$ aligned flow records reported above, the
partition sizes are those in Table~\ref{tab:chronological_split}.

\begin{table}[t]
    \centering
    \caption{Principal chronological flow partition.}
    \label{tab:chronological_split}
    \begin{tabular}{@{}lrrl@{}}
        \toprule
        Partition & Ratio & Number of flows & Temporal role \\
        \midrule
        Training & 70\% & 295,329 & Earliest flows \\
        Validation & 10\% & 42,189 & Intermediate flows \\
        Test & 20\% & 84,380 & Latest flows \\
        \midrule
        Total & 100\% & 421,898 & \\
        \bottomrule
    \end{tabular}
\end{table}

All packets belonging to one \datasetfield{flow_id} follow their parent flow,
so a flow cannot cross partitions. The preprocessing object is fitted only
with training-flow packets and training-flow statistics. Hyperparameter
selection, early stopping, and decision-threshold selection use validation
data; the test partition is reserved for final evaluation. This design is more
representative of deployment on later traffic than random row splitting and
reduces leakage caused by mixing packets or temporally adjacent copies of the
same behavior across partitions \cite{SommerPaxson2010,Arp2022DosDonts}.

The implementation also supports a chronological boundary search within a
specified tolerance when \datasetfield{stratify=true}. That optional mode
preserves contiguous time blocks but moves boundaries to reduce class-rate
differences. It is not active in the principal configuration described here;
if it is used in an experiment, the actual saved split counts and time ranges
must replace Table~\ref{tab:chronological_split}.

\subsection{Stage~1 Artifacts Supplied to Stage~2}
\label{subsec:stage2_artifacts}

After training, Stage~1 exports one learned embedding and one metadata record
per flow for each partition. The metadata include
\datasetfield{flow_id}, target label, split name,
\datasetfield{flow_start_timestamp_us}, source identifier, and destination
identifier. Stage~2 consumes the embedding as the content representation and
uses only the time and host identifiers to select eligible historical flows.
Keeping the original split label with every exported record prevents Stage~2
from accidentally constructing a context across train, validation, and test
boundaries. Labels are used as training/evaluation targets and are not included
among context features.

For reproducibility, each Stage~1 run also stores the selected feature names,
encoder categories, clipping bounds, scaler parameters, sequence policy,
random seed, ordered flow identifiers, class counts per partition, model
configuration, and checkpoint identifier. Stage~2 results are meaningful only
when they can be traced to the exact Stage~1 checkpoint and split manifest that
generated the embeddings.

\subsection{Validity, Privacy, and Reproducibility Considerations}
\label{subsec:dataset_limitations}

The resulting dataset is well suited to evaluating the proposed hierarchy on
a large, temporally ordered enterprise trace, but four limitations bound the
claims that can be made from it.

\begin{enumerate}[label=\textbf{L\arabic*:},leftmargin=*]
    \item \textbf{Alert-derived ground truth.} The labels inherit the coverage
    and errors of one Suricata ruleset. Negative records are more precisely
    described as ``no alert observed'' than as verified benign traffic.

    \item \textbf{Single-trace scope.} One company capture cannot by itself
    establish cross-organization, cross-sensor, or long-term generalization.
    External-PCAP evaluation is therefore a desirable extension.

    \item \textbf{Capture visibility.} The PCAP contains only traffic visible
    at the collection point. Packet loss, asymmetric routing, encrypted
    protocols, and incomplete flows can alter extracted statistics. The
    capture tool, interface location, capture interval, packet-loss report, and
    snap length should be documented when organizational policy permits.

    \item \textbf{Sensitive identifiers.} Raw addresses are retained locally
    to form Stage~2 relations. They must not appear in the thesis, shared CSVs,
    screenshots, or public metadata. A stable pseudonymization step should be
    applied before any release and verified not to change equality-based
    context membership.
\end{enumerate}

These constraints do not invalidate the experiment, but they determine its
interpretation. The study evaluates a reproducible model pipeline on a large
proprietary, weakly labelled capture. It does not claim that Suricata labels
constitute perfect attack truth, that all zero-day attacks are represented, or
that performance will transfer unchanged to every enterprise network. Clear
separation of raw fields, model predictors, context keys, and targets is used
throughout the remaining methodology and results sections to keep those claims
appropriately scoped.
